\section{Research question and verified gap}
AI-assisted equity valuation creates a strategic environment: an issuer chooses what to disclose, an evaluator converts that disclosure into a public valuation signal, and investors use the signal when allocating scarce capital. The issuer knows the evaluator's rule, so disclosure is potentially strategic. The deeper problem begins when the evaluator learns from later firm outcomes. If its own valuation affects how much capital a firm receives, and capital affects later performance, then part of the evaluator's training data is generated by the evaluator itself.

\textbf{Shared question.} \emph{When firms know how an AI evaluator scores their disclosures, can a history-aware evaluator and an auction-based capital-allocation mechanism reduce strategic disclosure while preserving informative pricing?}

We answer this through three connected questions. \textbf{Q1--strategic behavior:} how does a firm respond when it knows the evaluator and the evaluator's output affects capital? \textbf{Q2--mechanism and computation:} can evaluator learning separate genuine firm quality from performance caused by its own allocation? \textbf{Q3--behavior:} do investors remove predictable disclosure inflation from the AI signal, as the equilibrium benchmark assumes?

\textbf{Verified gap.} Costly-disclosure models study strategic reporting under a fixed evaluation rule~\cite{lacker1989,crocker2007}. IPO research studies uniform-price auctions and bookbuilding as information and allocation mechanisms~\cite{benveniste1989,sherman2005,degeorge2007}. Performativity work studies how deployed predictors change the environment they predict~\cite{perdomo2020}. These literatures do not jointly model a predictor that both affects capital allocation and later learns from outcomes partly caused by that allocation. Our bounded gap is therefore the feedback loop, not the generic claim that ``AI matters in finance.''

\section{Strategic model and solution concept}
There are sequential issuers $t=1,\ldots,N$, investors, and a mechanism designer. Issuer $t$ privately knows fundamental quality $\theta_t$, private value of capital $v_t$, and social marginal product of capital $\rho_t$. It chooses disclosure
\begin{equation}
d_t=\theta_t+m_t, \qquad c(m_t)=m_t^2/(2\gamma).
\end{equation}
The evaluator publishes
\begin{equation}
s_t=(1-\beta)\mu_0+\beta\,[a+b(d_t+\eta_t)],
\end{equation}
where $b$ is the learned disclosure sensitivity. Investors observe $s_t$ and private signals $x_{it}=\theta_t+\xi_{it}$, then bid. The allocation rule determines price $P_t$ and proceeds $k_t$. Finally,
\begin{equation}
y_t=\theta_t+\rho_t k_t+\varepsilon_t
\end{equation}
becomes part of the history used for later calibration.

The information structure is incomplete because $\theta_t,v_t,\rho_t$ are private types, and imperfect because manipulation $m_t$ is not directly observed. The game is sequential, so we use \textbf{Perfect Bayesian Equilibrium}~\cite{harsanyi1968}. Under the quadratic cost and single-crossing condition, the firm's best response is increasing in the evaluator's perceived sensitivity:
\begin{equation}
m_t^\star=\gamma v_t n\lambda_P[\lambda\beta b+(1-\lambda)\kappa_0].
\end{equation}
Thus revealing the evaluation rule changes behavior. The benchmark prediction is not that firms literally fool a fully rational evaluator: investors can anticipate average inflation. The economic cost can instead appear as resources spent to optimize for the machine.

\textbf{Challenge test.} If a history-aware evaluator learns a larger $b$ from outcomes it partly caused, then the best response above increases even when fundamental quality has not changed. This creates a falsifiable prediction: naive history should increase manipulation relative to a no-memory baseline.

\section{Social choice and design objective}
The social planner does not simply maximize AI prediction accuracy. The relevant stakeholders are issuers, institutional investors, retail investors, exchanges/regulators, and later issuers who inherit the evaluator's calibration. We therefore track three competing objectives:
\begin{enumerate}
\item \textbf{Informative pricing:} minimize error between market price and latent fundamental value.
\item \textbf{Allocative efficiency:} direct scarce capital toward firms with high social marginal product $\rho$.
\item \textbf{Manipulation and access:} reduce disclosure-management waste while preserving meaningful retail participation.
\end{enumerate}
These objectives can conflict. A mechanism that produces strong institutional information may reduce retail access; a transparent evaluator may improve accountability while also making the target easier to manipulate. The project therefore reports a trade-off rather than a single ``best'' rule.

The feedback loop also creates an intertemporal externality: one issuer's strategic disclosure can change the evaluator's future calibration for later issuers. A socially acceptable evaluator should therefore distinguish genuine information from performance induced by its own previous allocations.

\section{Mechanism design and auction comparison}
The scarce resource is 100 shares in an IPO-style offering. We compare two consequential allocation rules. \textbf{M1: uniform-price auction} sorts price--quantity bids and allocates until shares are exhausted; winners pay the lowest accepted price. \textbf{M2: bookbuilding} uses investor indications and underwriter discretion, with the current model assigning 85\% of allocation to institutions and applying an 8\% pricing discount. These are deliberately different mechanisms, not two labels for the same rule.

The key innovation is that the auction also creates variation in proceeds. Write the observed outcome as $y=\theta+\rho k+\varepsilon$. A naive learner that regresses outcome on disclosure obtains
\begin{equation}
\hat b_{naive}=b_{true}+\rho\frac{\operatorname{Cov}(k,d)}{\operatorname{Var}(d)}+R,
\end{equation}
where the middle term is the evaluator's own capital footprint and $R$ captures remaining heterogeneity. We therefore compare three evaluator modes: \textbf{no memory}, \textbf{naive history}, and \textbf{capital-purged history}. The last uses allocation variation as an instrument for $k$ and is gated by a first-stage strength check.

This makes mechanism design consequential. The question is not merely which auction raises more revenue. It is whether the allocation rule supplies enough credible variation to identify the feedback that the evaluator must remove before learning from outcomes. The analogy boundary is explicit: real bookbuilding also includes certification and relationship effects that our controlled model does not represent.

\section{Computational and behavioral test}
The computational design uses 200 synthetic issuers, 10 cohorts, seeds 42--46, 100 shares per offering, 15 institutional and 35 retail investors, and the same disclosure, signal, memory, and auction rules described above. The primary outcomes are mean manipulation $\bar m$, price RMSE, capital-allocation efficiency, feedback bias $\hat b-b_{true}$, manipulation cost, retail allocation share, and the first-stage $F$ statistic.

The current simulation reports, for example, $\hat b=1.42$ versus $b_{true}=1.02$ under naive memory and an increase in mean manipulation from 0.233 to 0.402. The current bookbuilding condition produces substantially stronger allocation-shock variation than the uniform-price condition. These are model outputs, not evidence about real issuers; all reported numbers must be regenerated from the final code and seed record before submission.

The behavioral artifact tests the equilibrium assumption that investors remove predictable disclosure inflation from the public AI signal. A competing explanation is anchoring or algorithm appreciation~\cite{logg2019}. Players therefore face the same strategic setting as the model and their implied weight on the public signal is compared with the Bayesian benchmark. Classroom/Hugging Face evidence is exploratory: it can challenge or qualify the behavioral assumption but cannot establish population-level investor behavior.

\textbf{Integrated contribution.} Game theory predicts strategic disclosure; social choice defines the collective objectives and trade-offs; mechanism design changes the allocation rule and the information available for evaluator learning; computation tests the feedback mechanism; and behavioral play challenges the rational-investor benchmark. The practical implication is testable: any AI system that both evaluates and influences capital should expose enough of its calibration path and allocation footprint for an external auditor to distinguish genuine performance from performance it helped create.
